\chapter{L'implementazione} % 5

Questo capitolo ha l'obiettivo di trattare l'implementazione del sistema progettato. Verranno quindi discusse le scelte implementative effettuate e quindi descritto l'implementazione dei singoli componenti. Tutto ciò verrà fatto con diversi livelli di dettaglio, andando più nel profondo per cose interessanti/difficili, più brevemente invece per quelle banali.

% È importante sottolineare come le scelte tecnologiche e le modalità di sviluppo qui presentate costituiscano una sola delle possibili soluzioni per la messa in atto del progetto.

\iffalse

L'architettura proposta:

\textbf{Natura statica dell'analisi}.
L'intero processo di analisi del campione dovrà essere basato esclusivamente su tecniche di analisi statica. Il sistema non dovrà quindi richiedere l'esecuzione del campione, né ricorrere a tecniche di analisi dinamica, debugging o osservazione diretta del comportamento durante l'esecuzione.
L'architettura basata su **Python MCP e Ghidra MCP con sole capacità statiche** rappresenta un approccio orientato al reverse engineering assistito da LLM, in cui i due componenti assumono ruoli complementari. Ghidra fornisce al modello una rappresentazione strutturata del programma analizzato, mettendo a disposizione informazioni quali funzioni, decompilato, riferimenti incrociati, grafi delle chiamate, stringhe e porzioni di memoria o byte selezionati. Python MCP agisce invece come ambiente di elaborazione generale, permettendo al modello di costruire ed eseguire strumenti personalizzati per trasformare, analizzare e interpretare le informazioni estratte da Ghidra.
Uno dei principali vantaggi di questa architettura è la flessibilità offerta da Python. Poiché il modello linguistico può generare codice Python durante l'analisi, è in grado di creare strumenti specifici per il problema affrontato senza richiedere un supporto predefinito. Ad esempio, quando viene individuata una struttura dati cifrata o una tabella di simboli offuscata, Ghidra può fornire i byte rilevanti mentre Python può essere utilizzato per implementare algoritmi di decodifica, testare ipotesi, confrontare risultati o ricostruire il formato dei dati. Questo rende il sistema particolarmente efficace in attività come l'analisi di configurazioni nascoste, il riconoscimento di pattern, la manipolazione di dati binari e la generazione automatica di strumenti di analisi.
L'utilizzo congiunto di Ghidra statico e Python permette inoltre di automatizzare attività complesse che richiederebbero normalmente numerosi passaggi manuali da parte di un analista. Ghidra consente di navigare la struttura del binario e identificare aree di interesse, mentre Python permette di effettuare elaborazioni più generali, come analisi statistiche, clustering, confronto tra campioni, generazione di report o integrazione con librerie esterne. In questo senso, Python MCP non sostituisce le capacità di reverse engineering di Ghidra, ma estende il sistema fornendo al modello la possibilità di creare procedure personalizzate sulla base delle informazioni recuperate.
Un confronto con un'architettura più completa, composta da **Python MCP e Ghidra MCP con capacità dinamiche di debugging ed esecuzione tramite P-code emulator**, evidenzia una differenza principalmente legata alla possibilità di osservare il comportamento runtime del programma. L'integrazione di strumenti dinamici permetterebbe al modello di eseguire porzioni di codice, osservare registri, memoria e flusso di esecuzione, verificando direttamente ipotesi formulate durante l'analisi statica. Questa capacità risulta particolarmente utile in presenza di codice offuscato, routine di unpacking, algoritmi di decrittazione con chiavi generate dinamicamente o comportamenti difficili da ricostruire esclusivamente dal codice.
Tuttavia, l'aggiunta di capacità dinamiche non rende superata l'architettura basata esclusivamente sull'analisi statica. La maggior parte delle attività preliminari di reverse engineering consiste infatti nella comprensione della struttura del programma, nell'identificazione delle funzioni rilevanti e nella ricostruzione delle relazioni tra componenti. In questi scenari, l'analisi statica combinata con un ambiente di scripting flessibile può essere più efficiente rispetto all'esecuzione diretta del codice, soprattutto quando è necessario analizzare numerosi campioni o creare pipeline automatizzate.
In conclusione, le due architetture possono essere considerate complementari. L'approccio **Python MCP + Ghidra statico** fornisce un ambiente versatile per l'analisi, la trasformazione dei dati e la costruzione automatica di strumenti, risultando sufficiente per molte attività di reverse engineering assistito. L'aggiunta delle capacità dinamiche di Ghidra introduce invece un ulteriore livello di verifica sperimentale, permettendo al modello di passare dalla deduzione del comportamento alla sua osservazione diretta. Un sistema completo potrebbe quindi utilizzare l'analisi statica come fase principale di comprensione e ricorrere all'esecuzione dinamica solo quando necessario per validare ipotesi o analizzare comportamenti non osservabili staticamente.
\fi

\iffalse
Ma c'è anche un aspetto spesso sottovalutato: dare a un LLM un debugger completo non significa automaticamente ottenere un agente migliore. Aggiungere capacità dinamiche aumenta anche la complessità:
deve decidere cosa eseguire;
deve preparare un ambiente corretto;
deve interpretare trace molto grandi;
deve evitare di perdersi in dettagli irrilevanti;
deve gestire casi in cui l'esecuzione non è rappresentativa.
Un agente con molti strumenti non è necessariamente più efficace di un agente con gli strumenti giusti.
La tua scelta ha anche un vantaggio progettuale: hai separato bene i ruoli.
Ghidra statico: fonte autorevole della conoscenza sul binario.
Python: ambiente dove il LLM può creare procedure e ragionamenti personalizzati.
Questa combinazione è già una base molto forte. Se in futuro volessi estenderla, il debugging/P-code sarebbe probabilmente un terzo livello, non una sostituzione:
\fi

\section{Ambiente di sviluppo} % 5.1

L'implementazione del sistema è stata realizzata in Python 3.13 \footnote{\url{https://www.python.org}}, mentre per la gestione dell'ambiente di sviluppo e delle dipendenze del progetto è stato utilizzato Conda \footnote{\url{https://docs.conda.io/projects/conda/en/latest/user-guide/tasks/index.html}}.

Il motivo per cui la scelta del linguaggio di implementazione è ricaduta su Python è legato principalmente all'ampio ecosistema di librerie disponibili per le funzionalità richieste dal sistema. Tra queste rientrano, in particolare, librerie dedicate all'analisi del malware, all'integrazione con LLM e agenti, alle interazioni basate sul protocollo MCP e all'interazione con gli ambienti di esecuzione. La disponibilità di strumenti già consolidati ha pertanto consentito di concentrare lo sviluppo sulla logica specifica dell'architettura proposta, riducendo la necessità di implementare direttamente funzionalità di supporto.

% https://www.python.org/downloads/release/python-3130/
% https://docs.conda.io/projects/conda/en/latest/user-guide/tasks/manage-python.html

\section{Orchestrazione / Esecuzione} % 5.2

\subsection{Classifier} % 5.2.1

\subsubsection{Virus Total (?)}

API: malevolo oppure non malevolo e ipotesi famiglia

\subsubsection{YARA Rules (?)}

Conferma famiglia.

\subsection{SRE Agent} % 5.2.2

\begin{itemize}
    \item import in SRE Framework e analisi automatica (fun ghidra)
    \item determina con una funzione se stripped (fun ghidra)
    \item eventuale injection knowledge base
    \item doppio ciclo di renaming start/stop
    \item save (fun ghidra)
    \item doppio ciclo di reversing con due report
\end{itemize}

contatto con SRE framework e sandbox.

\subsubsection{Knowledge Base}

\subsubsection{Selezione Tools Analisi Statica} % 5.3.4

Per questa implementazione è stato scelto solo un sottoinsieme di tool riguardanti analisi statica del codice.
tre ragioni:
- semplicità, cioè riduzione delle variabili in gioco e più controllo sulle azioni (eventualmente estendibile)
- evitare di esporre all'esecuzione reale
- the less mcp the better (LLM)

I tool scelti rientrano nelle seguenti categorie:
Tools YAML.
[motivazioni parte 1]

\subsubsection{Selezione Tools Analisi Dinamica}

Niente accesso al file. niente accesso a internet.

[motivazioni parte 2] (no malware access)

\subsection{Code Agent} % 5.2.3

\subsubsection{Context Injection}

\section{SRE Framework} % 5.3

Il framework di SRE adottato è Ghidra, sviluppato e mantenuto dal Research Directorate della National Security Agency (NSA) e distribuito come progetto open source secondo la Apache License 2.0.

L'ambiente Ghidra, \textit{platform independent}, offre una suite ricca di strumenti e funzionalità per l'analisi binaria del codice compilato, tra cui disassembly, decompiling, graphing, debugging e scripting. Il framework supporta inoltre un'ampia varietà di \textit{instruction set} relativi a differenti architetture di processore e numerosi formati eseguibili. Le attività di analisi possono essere svolte sia in modalità interattiva, attraverso una \textit{Graphical User Interface} (GUI), sia in modalità automatizzata, rendendo Ghidra adatto all'integrazione all'interno di processi di analisi più complessi \cite{nance2026ghidrabook}.

% Cite NSA

\subsection{Ghidra} % 5.3.1

\begin{itemize}
    \item Headless (no GUI)
    \item Docker Image
    \item \footnote{\url{https://github.com/nationalsecurityagency/ghidra}}
\end{itemize}

\subsection{MCP Plugin} % 5.3.2 (ghidra-mcp)

Per integrare le funzionalità di analisi offerte da Ghidra con il \textsc{SRE Agent} in conformità con il protocollo MCP è stato utilizzato il software open source \texttt{ghidra-mcp}, disponibile sull'omonimo repository GitHub \footnote{\url{https://github.com/bethington/ghidra-mcp}} e distribuito secondo la licenza Apache License 2.0.

Il software fornisce un'estensione per Ghidra e un insieme di componenti di supporto che permettono di rendere accessibili le funzionalità dell'ambiente di analisi attraverso il protocollo MCP, esponendo complessivamente più di duecento tool. Si tratta di una soluzione in continua evoluzione che supporta l'integrazione con Ghidra Server, la gestione di più progetti e file, l'esecuzione sia in modalità GUI sia in modalità headless e il deployment mediante Docker.

% meccanismi di gestione errori e affidabilità

L'insieme delle funzionalità messe a disposizione è particolarmente ampio e comprende, innanzitutto, le principali operazioni di analisi binaria, relative a funzioni, data flow, strutture dati, stringhe, operazioni di importazione ed esportazione e accesso alla memoria. A queste si aggiungono funzionalità più avanzate, tra cui, nell'ambito dell'analisi dinamica, l'emulazione del P-code e l'integrazione con il debugger, oltre a strumenti per la gestione dei progetti, del controllo di versione, degli script di Ghidra. Sono inoltre supportate l'esecuzione batch delle operazioni e l'analisi di più programmi contemporaneamente.

% Analysis control

\subsection{Ghidra MCP Server} % 5.3.3 (bridge)

Il server MCP del software \texttt{ghidra-mcp} è implementato mediante il modulo \texttt{bridge\_mcp\_ghidra}, sviluppato in Python utilizzando la libreria FastMCP \footnote{\url{https://gofastmcp.com/}}. Il modulo costituisce il livello intermedio tra i client MCP e l'estensione di Ghidra e supporta sia il trasporto \textit{stdio}, sia il trasporto \textit{Streamable HTTP}.

Una caratteristica interessante del modulo è l'organizzazione dei tool in gruppi semantici, che permette di gestire in maniera selettiva l'insieme delle funzionalità esposte al LLM. In modalità standard, il server rende disponibili tutti i gruppi di tool; tuttavia, considerata l'ampiezza del catalogo, il software può essere eseguito anche in modalità \textit{lazy}, caricando inizialmente solo i gruppi di tool principali. Gli altri gruppi possono essere individuati e successivamente caricati o rimossi durante l'esecuzione da parte del LLM, in base alle necessità.

% \texttt{search_tools}, \texttt{list_tool_groups}, \texttt{load_tool_group}, \texttt{unload_tool_group} e \texttt{check_tools}

\section{Ambiente Sandbox} % 5.4

La sandbox di SRE viene realizzata mediante un ambiente di esecuzione containerizzato, utilizzando Docker per la creazione, la configurazione e la gestione dei container.

La containerizzazione è una tecnica di isolamento a livello di sistema operativo che consente di eseguire un'applicazione e le relative dipendenze all'interno di un ambiente isolato. Rispetto alle macchine virtuali, che richiedono l'esecuzione di un intero sistema operativo guest, i container rappresentano un approccio più leggero per ottenere ambienti di esecuzione separati \cite{docker2026container}.

Nel contesto della \textsc{SRE Sandbox}, Docker fornisce i meccanismi necessari alla creazione e alla configurazione dei container, consentendo di controllarne il filesystem, l'accesso alla rete e le risorse computazionali disponibili. Queste caratteristiche permettono di configurare l'ambiente in modo coerente con i requisiti di isolamento e sicurezza individuati durante la fase di progettazione.

\subsection{Configurazione} % 5.4.1 del Dockerfile

La configurazione dell'ambiente \textsc{SRE Sandbox} è definita mediante un \textit{Dockerfile}, che specifica l'immagine di base, le dipendenze e le impostazioni necessarie al funzionamento del container. L'immagine utilizzata è basata su Python 3.13 e Debian Slim \footnote{\url{https://hub.docker.com/_/python}}, così da fornire un ambiente contenuto, ma sufficientemente completo da permettere al LLM di creare e modificare file, eseguire codice Python e impartire comandi tramite la shell.

A questo scopo, durante la fase di costruzione vengono installati gli strumenti di sistema e le librerie Python necessarie alle attività previste nella sandbox, tra cui quelle dedicate alle funzionalità crittografiche, come \texttt{cryptography} \footnote{\url{https://cryptography.io}} e \texttt{pycryptodome} \footnote{\url{https://www.pycryptodome.org/}}. Viene inoltre configurata una directory dedicata alla sandbox e viene creato un utente senza privilegi di root, utilizzato per l'esecuzione dei processi all'interno del container anziché l'utente root.

L'immagine così definita viene utilizzata per creare il container della sandbox, che implementa ulteriori misure di isolamento. In particolare, l'accesso alla rete viene completamente disabilitato, limitando l'interazione del codice con risorse esterne e confinando le operazioni alle sole funzionalità e dipendenze disponibili all'interno dell'ambiente. Per consentire lo scambio controllato dei \textsc{Sandbox Files}, una directory dell'host viene montata all'interno del container con permessi di lettura e scrittura. Infine, vengono imposti limiti alle risorse computazionali disponibili, in termini di memoria e capacità di elaborazione.

\subsection{Sandbox MCP Server} % 5.4.2

Il server MCP della \textsc{Sandbox} ha il compito di interagire con l'ambiente di esecuzione ed esporne le funzionalità attraverso tool MCP. Per la sua implementazione viene utilizzata la libreria Python \textit{FastMCP} \footnote{\url{https://gofastmcp.com/}}, mentre l'interazione con il container viene gestita attraverso la libreria \textit{Docker SDK for Python} \footnote{\url{https://docker-py.readthedocs.io}}.

All'inizializzazione del server viene verificata la disponibilità del container della sandbox e, qualora questo non risulti in esecuzione, viene avviato automaticamente. Il server espone quindi un insieme di tool organizzati in due principali categorie:

\begin{description}
    \item[File System] I tool \texttt{list\_files}, \texttt{read\_file} e \texttt{write\_file} permettono di consultare e modificare i file presenti nell'ambiente della sandbox. Ogni operazione prevede inoltre una gestione degli errori, così da fornire un'informazione esplicita in caso di esecuzione non riuscita.

    \item[Execution] Il tool \texttt{execute\_python\_file} consente di eseguire file Python all'interno del container, imponendo un limite temporale all'esecuzione. Il tool \texttt{execute\_command} permette invece di eseguire comandi della shell. Prima dell'esecuzione viene effettuato un controllo di sicurezza sui comandi richiesti per prevenire operazioni potenzialmente pericolose, come tentativi di privilege escalation, o apertura di nuove shell. In entrambi i casi, il risultato dell'esecuzione viene restituito in forma strutturata, includendo l'\textit{exit code}, \texttt{stdout} e \texttt{stderr}.

    % Controllo della shell syntax
\end{description}

\section{MCP Client} % 5.5

Per i \textsc{Sandbox MCP Client} e \textsc{Ghidra MCP Client} è stata definita una classe base comune che si occupa di instaurare e mantenere la connessione con il rispettivo server mediante il metodo di trasporto \textit{stdio}. La classe gestisce il ciclo di vita della connessione attraverso un contesto asincrono, inizializzando una ClientSession (python mcp sdk) e fornendo le due funzioni di base list tools e call tool.

\subsection{LLM-MCP Client} % 5.5.1 Tool Parsing

La classe viene quindi estesa, introducendo quindi il livello di integrazione con il LLM. In particolare gli strumenti esposti dal server MCP vengono convertiti nel formato atteso dalle API di OpenAI, quello di tool/function generation.

Infine, un'altra classe (???) fornisce anche un meccanismo per la registrazione di un client MCP: consente di selezionare, attraverso i nomi un sottoinsieme di strumenti da esporre e di applicare un prefisso (id del MCP) ai nomi dei tool.

Nel complesso, queste componenti costituiscono un livello di astrazione che semplifica la gestione dei diversi client MCP e il loro collegamento con il LLM per il tool calling del \textsc{SRE Agent}.

\section{LLM Client} % 5.6

Questo è il componente responsabile dell'interazione tra gli agenti, il LLM e gli strumenti esposti dai client MCP.

All'inizializzazione, gli strumenti disponibili vengono raccolti (dai MCP client passati) e associati al rispettivo client mediante una struttura di risoluzione, permettendo in seguito di individuare il server MCP responsabile di ciascuna richiesta generata dal LLM.

Ci sono le funzionalità di impostazione dei parametri di utilizzo del LLM, in particolare temperatura, \dots.

Vengono quindi messi a disposizione dei metodi per risposte singole (chat completion secondo lo standard OpenAI), con e senza l'utilizzo dei tool. Poi viene esposto il metodo per il loop dell'agente. Il loop viene implementato in modalità lossless, senza compattazione di contesto (in quanto non è normale andare oltre per le task decise in questa implementazione).

Nel tool dell'agente ovviamente c'è il dispatching dei tool chiamati, eseguendoli sul rispettivo MCP, aggiungendo al contesto le risposte formattate ricevute dagli ambienti.

\subsection{Logging} % 5.6.1

per l'agente è necessario implementare le funzionalità di logging. a questo scopo viene fatto su due punti di vista (?) / livelli diversi:

- uno completo, con il json sempre riferimento a standard OpenAI di intere conversazioni
- uno orientato alla lettura umana: prompt iniziale, numero di step, trascrittura di ragionamento e assistant response, con tool chiamati (ma senza le risposte).

% \subsection{Tool Calls Handling} % 5.6.2

\section{MessageSchema} % 5.7

Come individuato nella fase di analisi dei requisiti e progettazione, è necessario definire uno schema comune per i messaggi restituiti dal \texttt{Python Module}. A questo scopo è perfetto utilizzare uno JSON schema. 

JSON Schema è un linguaggio dichiarativo progettato per descrivere la struttura e i vincoli dei dati rappresentati in formato JSON. Attraverso un insieme di keyword consente di specificare la tipologia dei dati, la struttura degli oggetti e delle proprietà che li compongono \cite{jsonschema2020json}.

L'uso di JSON schema è ideale per due ragioni. La prima è che permette di verificare la conformità di un documento (istanza?) JSON a uno schema mediante un apposito \textit{validator}, confrontando l'istanza con le regole definite nello schema, determinando se i vincoli specificati sono rispettati. Oltre allo scopo di validazione, JSON schema può anche essere utilizzato come forma di documentazione, fornendo una descrizione di struttura e vincoli dei dati attraverso una rappresentazione comprensibile sia dall'umano che dal LLM.

Da qualche parte bisogna affermare la compatibilità con i dizionari python.

\subsection{Message Classification} % 5.7.1

Prima di procedere all'implementazione dello schema, conviene sfruttare l'espressività (???) di JSON schema per poter distinguere (semanticamente?) i messaggi catturati dall'infrastruttura di C\&C in diverse tipologie, facilitando (per l'umano capire quali messaggi vengono scambiati) e adottando ricorsivamente schemi diversi in base al tipo.

Per questa implementazione particolare si è deciso di distinguere tra i seguenti tipi di messaggi:

\begin{description}
    \item[ACK] per i messaggi di richiesta di invio acknowledgment a un precedente messaggio richiesti dal server C\&C.
    \item[Attack] per i messaggi che contengono istruzioni di attacco.
    \item[Heartbeat] per i messaggi di keepalive del traffico.
    \item[Command] per i messaggi di servizio o generici, che quindi non rientrano nelle categorie precedenti oppure che non si è convinti di classificare più specificamente.
\end{description}

\subsection{JSON Schema} % 5.7.2

Lo schema realizzato è riportato tra le appendici. è conforme al seguente standard \footnote{\url{https://json-schema.org/draft/2020-12/json-schema-core.html}}

Al primo livello dello schema si decide di salvare timestamp temporale di ricezione del messaggio e indirizzo IP sorgente del messaggio. Quindi c'è un campo per identificare il tipo del messaggio tra quelli precedenti e poi uno schema annidato con le informazioni specifiche per il messaggio.

Descrizione specifiche dei messaggi: ACK ha almeno uno status, il heartbeat ha un numero di sequenza (lato-server) e uptime opzionale, mentre command riporta il payload raw e una stringa di tipo del comando (o null) se si ha un'identificazione specifica. Tutti i messaggi hanno una nota string/null per informazioni aggiuntive.

Il tipo attack è il più interessante: codice identificativo dell'attacco, nome specifico dell'attacco e protocollo di rete coinvolto, durata dell'attacco, numero e valore degli argomenti passati all'attacco (argc, argv), numero e indirizzi ip dei bersagli dell'attacco.

\subsubsection{IP Schema}

Nello schema c' anche un campo array per l'info degli IP coinvolti dagli attacchi, eventualmente vuoto.

Questo permette di avere informazioni specifiche sui bersagli coinvolti dagli attacchi: tra queste informazioni ci sono quelle geografiche (da continente fino a città), di localizzazione (latitudine, longitudine, timezone), sulla rete a cui appartiene (CIDR, ASN e name, domain), info DNS e spazio per eventuali aggiunte.

\subsubsection{IP Lookup}

La responsabilità di "compilare" queste informazioni, a partire dagli IP dell'attacco, è del Python Module. Tuttavia questo deve solo invocare un modulo fisso.

-> riscrivi per chiarezza. il modulo è fisso e realizzato dall'utente, passato direttamente al Module.

nella pratica per implementare ci si può affidare a database locali da cui estrarre le informazioni, insieme a librerie ipaddress e dnsPython.
% https://ipinfo.io/data/ipinfo_lite.mmdb?token=b2626200bf26b3
% https://download.maxmind.com/geoip/databases/GeoLite2-ASN/download?suffix=tar.gz

\section{Code Templates} % 5.8

% \subsection{Aggiornamento BotNetInterface} % 5.8.0

Come anticipato nella progettazione, conviene separare le responsabilità, per cui si utilizzeranno due file distinti, partendo comunque dal BotNet Module che rimane il main.

Visto che nel modulo verranno implementate funzioni di logging e c'è la dipendenza dalla classe centralizzata IP lookup, a botnet interface si aggiunge un costruttore che permette di passare questi oggetti.

Per la classificazione dei messaggi si usa una ENUM python (MessageClassification).

\subsection{BotNet Module} % 5.8.1

deve implementare  secondo i contratti definiti da BotNetInterface, le funzioni del modulo main principale, andando un po' più nello specifico, secondo quello che ci si aspetta dal modulo finale.

\begin{description}
    \item[Attributi/Costanti di classe] hostname o host ip, sequenze di handshake, ack, heartbeat ==> tutti i parametri (se) fissi del modulo (aggiungi tu)
    \item[Costruttutore] creazione di quanto necessario, del socket, del buffer di ricezione dei messaggi, del logger e del modulo parser dei messaggi (classe).
    \item[Connect] stabilire praticamente una connessione all'endpoint C\&C del socket e implementazione tramite funzione authentication dedicata della sequenza di inizializzazione del protocollo, così da potersi registrarsi come se fosse il dispositivo di una vittima al server di C\%C.
    \item[Sniff] rimane in ascolto indefinito dei comandi ricevuti dal C\&C e poi c'è il metodo di heartbeat regolare. Quando riceve qualcosa applica un eventuale pre-processing allo stream (preprocess\_stream) ed estende il buffer; dal buffer si estraggono singoli pacchetti (extractPackets) e si esegue il decoding o decifratura (decodePacket) se prevista dal protocollo; quindi dai pacchetti si estraggono i messaggi (extractMessages); per ciascun messaggio si classifica il tipo (getmessagetype, secondo i principi introdotti prima) e, se il protocollo lo richiede si adotta un comportamento specifico (ad esempio ACK ai messaggi ricevuti). Il parsing del messaggio è demandato al modulo MessageParser (funzione parseMessage), quindi il dizionario viene ritornato.
\end{description}

tutto viene implementato tramite "pass" python oppure return/packet degli oggetti senza modifiche (immutati). quello che si trasmette è che l'implementazione dovrà:

\begin{itemize}
    \item se il protocollo implementa quella cosa, modifica e replicala (OVERRIDE)
    \item se il protocollo non la implementa, lascia tutto com'è, visto che non modifica niente e passa avanti.
\end{itemize}

Affinché poi vengano passati al Code Agent, tutte le funzioni adottano commenti """""" per definire contratto/comportamento dei singoli funzioni (nei prompt quindi basterà più una descrizione di insieme e non specifica). durante la classe vengono anche mostrate informazioni riguardanti le informazioni/azioni per cui ci si aspetta il logging.

\subsection{MessageParser Module} % 5.8.2

deve esporre un metodo \texttt{parse\_message} che possa essere richiamato dal botnet module. questo metodo può quindi direttamente inserire le informazioni comuni a tutti i messaggi, successivamente viene eseguito il dispatch verso i metodi specifici (parse ack, command, heartbeat e attack). Nel caso di attack in più è responsabile di riempire le informazioni ipI  nfo utilizzando la rispettivamente classe.

Implementa la logica secondo cui un errore di parsing dei metodi specifici fa ricadere su un messaggio generico command. alla fine il metodo deve restituire un dizionario python pronto a essere validato contro il JSON schema.

Per semi implementazioni e commenti valgono stessi principi di botnet module.

\section{Prompt Archive} % 5.9 (Engineering)

\subsection{Renaming} % 5.9.1

\subsubsection{Start}

\subsubsection{Continue}

\subsection{Reversing} % 5.9.2

\subsubsection{Pt. 1}

\subsubsection{Pt. 2}

\subsection{Coding} % 5.9.3

\subsubsection{Botnet}

\subsubsection{Parser}

\subsubsection{YAML}
